\chapter{Análisis de Resultados}

La ejecución de los procedimientos de prueba sobre la RDB de \colorbox{gray!20}{\textcolor{lineText}{\texttt{import\_tech}}} \supercite{suarez2010} generó una serie de bloqueos nativos controlados por el motor MySQL, los cuales evalúan y validan la solidez del modelo físico implementado para Import Tech S.A.S. Al contrastar las sentencias DML anómalas con las respuestas del sistema, se obtuvieron los siguientes resultados definitivos:

\begin{itemize}
    \item \textbf{Validación de Integridad Referencial:} El sistema arrojó el \textbf{Error 1452} al intentar registrar transacciones con dependencias inexistentes (un detalle de venta sin producto y una venta vinculada a un cliente no registrado). Este resultado demuestra que las restricciones \colorbox{gray!20}{\textcolor{lineText}{\texttt{FOREIGN KEY}}} actúan correctamente, impidiendo la creación de datos huérfanos y garantizando la coherencia de las dependencias comerciales.
    
    \item \textbf{Validación de Integridad de Entidad:} La aparición del \textbf{Error 1062} ante los intentos de duplicar la llave primaria en la tabla \colorbox{gray!20}{\textcolor{lineText}{\texttt{categoria}}} y los atributos únicos (\colorbox{gray!20}{\textcolor{lineText}{\texttt{NIT}}} y \colorbox{gray!20}{\textcolor{lineText}{\texttt{NIP}}}) en la tabla \colorbox{gray!20}{\textcolor{lineText}{\texttt{proveedor}}}, certifica que el sistema respeta el principio de unicidad, imposibilitando la duplicidad de registros en el catálogo.

    \item \textbf{Validación de Integridad de Dominio:} El bloqueo de la transacción evidenciado mediante el \textbf{Error 1048} al omitir un campo obligatorio, confirma que las cláusulas \colorbox{gray!20}{\textcolor{lineText}{\texttt{NOT NULL}}} son efectivas. El motor exige estrictamente la completitud de la información crítica (como asociar la categoría a un producto) antes de procesar el guardado.

    \item \textbf{Validación de Restricción de Eliminación:} Al ejecutar operaciones \colorbox{gray!20}{\textcolor{lineText}{\texttt{DELETE}}} sobre categorías y proveedores que ya poseían productos asociados, el motor emitió el \textbf{Error 1451}. Este bloqueo es vital, ya que confirma que la base de datos protege de manera proactiva los registros padre, evitando que los productos del inventario queden sin una referencia hacia su distribuidor o clasificación.
\end{itemize}